\documentclass[11pt]{article}
\usepackage[T1]{fontenc}
\usepackage[utf8]{inputenc}
\usepackage{amsmath,amssymb}
\usepackage[margin=1in]{geometry}
\usepackage{url}
\setlength{\emergencystretch}{3em}
\title{Technical Review of Minimum-Payment Expense Settlement}
\author{Review of \texttt{docs/methodology.tex}}
\date{14 November 2026}
\begin{document}
\maketitle

\section{Assessment and scope}

The central characterization of the optimum as $n-c$ is correct, and the
three-way distinction between exact optimization, subset-packing heuristics,
and polynomial greedy settlement is useful. However, the manuscript needs
major revision before it supports its claims of proved correctness. There
are false statements, an invalid lower-bound argument, incorrect executable
pseudocode, and several contradictions in the complexity analysis.

This review compares the manuscript with the current
\texttt{settle\_expenses.py} and inspects
\path{tests/test_heuristic_vs_exact.py}. It is a mathematical and source
review, not an exhaustive execution audit or a performance benchmark.
The suggested corrections below distinguish errors in the exposition from
actual differences between the exposition and implementation.

\section{Major mathematical corrections}

\subsection{The $n-1$ theorem has the wrong quantifier}

Section 3 states that \emph{any feasible settlement} uses at most $n-1$
payments. Feasibility imposes no such upper bound: a debt of two cents can
be paid in two one-cent transfers, although $n=2$. Splitting a payment can
increase the count further whenever the amount permits it. The intended
claim is that \emph{there exists} a feasible settlement with at most $n-1$
payments, for $n>0$, and hence the optimum has that upper bound.

The proof should use a transfer of
$\min(d_i,c_j)$ between a debtor and a creditor. Sending the debtor's
entire debt to an arbitrary creditor can overpay that creditor and does not
describe the stated debtor--creditor algorithm. With the minimum rule,
each payment clears at least one person, and the final payment clears two.
This proves the bound. Treat $n=0$ separately: the optimum is zero,
not $n-1=-1$.

\subsection{Replace the proof of the $n-c$ characterization}

The lower-bound proof in Section 3 invokes an upper-bound construction
to claim a lower bound. Its displayed expression
$\sum_j(k_j-1)=n-\sum_j k_j$ is also incorrect: the left side equals
$\sum_j k_j-c$. Combining an upper bound on the remainder with these
claims cannot establish the theorem.

A complete replacement is as follows. Define $c$ as the maximum number
of \emph{nonempty} blocks in a zero-sum partition of $S$. Consider any
feasible settlement, and form its underlying undirected multigraph, with
one edge per payment. Summing the balance equations within a connected
component cancels all internal transfers. Each component therefore has
zero total initial balance. If there are $q$ components, they form a
zero-sum partition, so $q\leq c$. A connected component with $k$ vertices
needs at least $k-1$ edges. Consequently,
\[
 |P|\geq\sum_{j=1}^{q}(k_j-1)=n-q\geq n-c.
\]
Conversely, settle each block in a maximum zero-sum partition using at
most $k_j-1$ payments. This constructs a feasible settlement with at most
$n-c$ payments. The lower bound makes this count exact.

This proof assumes unrestricted transfers among the modeled participants.
State that assumption explicitly; restrictions on who may pay whom would
change the problem. If zero-balance intermediaries are permitted, explain
why they cannot improve the bound, or explicitly restrict the settlement
vertex set to $S$.

For nonempty zero-sum $S$, a maximum disjoint collection automatically
covers $S$: any nonempty remainder also sums to zero and could be added.
Thus $c\geq1$ and $\min(n-c,n-1)=n-c$. The repeated minimum expression
and the hypothetical case $c<1$ are unnecessary for $n>0$.
Define $c=0$ for the empty instance.

\subsection{Exactly $k-1$ payments requires an indecomposable block}

An arbitrary zero-sum group does not require exactly $k-1$ greedy payments.
For balances $(1,-1,2,-2)$, largest-to-largest matching uses two payments,
whereas $k-1=3$. In general the construction gives \emph{at most} $k-1$.

The stronger equality does hold for a block of a maximum-cardinality
zero-sum partition. Such a block has no proper nonempty zero-sum subset:
otherwise it could be split to increase the number of blocks. Any feasible
settlement within it must therefore be connected, and needs at least
$k-1$ payments. Supply this reasoning wherever equality is used.

\section{A faster exact algorithm should be discussed}

The manuscript's $O(3^n)$ algorithm is not the strongest available exact
subset DP. There is an $O(n2^n)$-time, $O(2^n)$-space formulation that
considers deleting one person rather than enumerating submasks. This is
a material omission from the comparison and the discussion of practical
thresholds.

Let $s(M)=\sum_{i\in M}b_i$, and let $f(M)$ be the maximum number of
nonempty zero-sum prefixes among all orderings of the people in $M$.
Then
\[
 f(\varnothing)=0,\qquad
 f(M)=\max_{i\in M}f(M\setminus\{i\})+
             \mathbf{1}_{\{s(M)=0\}}\quad(M\ne\varnothing).
\]
The recurrence follows by choosing the last person in the ordering:
earlier prefixes belong to the ordering of the predecessor mask, and
the full prefix contributes precisely when its sum is zero.

On the full zero-sum set $S$, consecutive zero-sum prefixes delimit
disjoint zero-sum blocks covering $S$. Conversely, concatenate the
blocks of any zero-sum partition in any internal order; each block
boundary is a zero-sum prefix. Hence $f(S)$ is exactly the maximum
partition count $c$, and the minimum payment count is $n-f(S)$.
This interpretation also defines the DP correctly on intermediate masks
whose total balance is not zero.

Precompute sums in $O(2^n)$ time by removing a lowest set bit. There
are exactly $n2^{n-1}$ single-person transitions across all masks for
$n\geq1$. Store a maximizing predecessor for each mask, recover the
ordering, and cut it wherever the running balance returns to zero.
Settling the recovered blocks then produces a minimum-payment result.
The empty instance returns zero directly.

For comparison, at $n=18$ this DP has $2{,}359{,}296$ transitions,
versus $193{,}710{,}244$ candidate checks in the manuscript's submask
DP. These are operation counts, not measured runtime predictions.
The exponential memory requirement still matters, and any new default
threshold should be benchmarked.

An equivalent minimum-cost formulation is documented in the official
AtCoder editorial for ABC432 F, \emph{Candy Redistribution}, under
``Solution in O(N * 2\^{}N)'':
\begin{center}
\texttt{https://atcoder.jp/contests/abc432/editorial/14584?lang=en}
\end{center}
The prefix-count derivation above supplies the connection to this
manuscript's $n-c$ characterization. The original $O(3^n)$ algorithm
may still be described accurately as the current implementation, but
should not be presented as the only exact alternative.

\section{Exact dynamic programming and reconstruction}

\subsection{The recurrence is sound, but its proof needs to be stated}

The two options in Section 4 correctly cover the possibilities that the
lowest-bit person is unused or belongs to a selected zero-sum block.
Use induction on the mask to prove that each option is feasible and that
every optimal collection is represented by one of them. All successor
masks are smaller, so the numerical evaluation order is valid.

Exclude the empty mask explicitly from \texttt{zero\_sum\_masks}, as the
implementation does. The claim that every partition is considered
``exactly once'' is stronger than necessary; coverage of every feasible
collection suffices. The final optimality proof cites Lemma 1 as proving
the DP optimum, but that lemma addresses reconstruction only. Introduce
and reference a separate recurrence-correctness result.

\subsection{A zero parent does not generally mean no subsets remain}

For the stated DP semantics, \texttt{parent[m] = 0} records a skip of
the lowest bit. It does not imply $\texttt{dp[m]}=0$. For example, take
balances $(2,1,-1)$ in bit order. At the full mask, the best collection
is the last two people, but the parent is zero because the lowest-bit
person is skipped. The reconstruction pseudocode immediately returns
an empty collection. This refutes the general reconstruction argument
given in the manuscript.

There is an important qualification: this example is not a zero-sum full
input, and it does not show that the current public exact method fails.
For a nonempty zero-sum mask $m$, Option A leaves the lowest-bit person
unused. Any collection obtained that way leaves a nonempty zero-sum
remainder in $m$, which can be added as another block. Hence Option B
strictly improves on Option A. A maximizing parent is therefore nonzero,
and removing it preserves a zero-sum remainder. By induction, reconstruction
from the full zero-sum mask reaches zero and extracts an optimal partition.

Either restrict and prove the reconstruction lemma using this invariant,
or make reconstruction valid for arbitrary masks by replacing the break
with a skip:
\begin{verbatim}
if sub == 0:
    m ^= m & -m
    continue
\end{verbatim}
The latter is a proposed generalization, not the current implementation.

\subsection{Remove the impossible optimal-remainder case}

Section 4 concludes both $|P|\leq n-c-1$ when a remainder exists and
$|P|\geq n-c$. These cannot both hold. With a correctly reconstructed
maximum collection on a zero-sum input, no remainder exists. Explain this
and state $|P|=n-c$ directly. The implementation's defensive remainder
branch can be described separately. It settles only the remaining people,
not the entire original balance set as the manuscript currently suggests.

\section{Greedy algorithm and implementation fidelity}

\subsection{The debtor heap sign is wrong in the manuscript}

For $b_i<0$, the expression $-b_i$ is positive. Thus the displayed debtor
heap selects the smallest debt first, and the pseudocode computes a
negative payment through $\min(-k_d,-k_c)$. With balances $(-5,+5)$,
it yields $\min(-5,5)=-5$. This violates the basic payment definition.

The actual implementation correctly uses the following keys:
\begin{verbatim}
debtors   = [(b_i, i)  for i if b_i < 0]
creditors = [(-b_i, i) for i if b_i > 0]
\end{verbatim}
Both keys are negative magnitudes. Include explicit residual updates:
after paying $a$, reinsert $(k_d+a,d)$ or $(k_c+a,c)$ only when the
corresponding key remains negative.

\subsection{Use the right progress invariant and operation count}

The signed sum of remaining balances stays zero, so it cannot establish
termination by decreasing. Instead, the total remaining debt decreases
by a positive integer amount on each iteration. More strongly, at least
one nonzero balance is eliminated per payment, giving at most $n-1$
iterations for $n>0$.

At termination \emph{both} heaps are empty: conservation of equal total
debt and credit rules out exactly one nonempty heap. Each iteration pops
two entries and reinserts at most one. The manuscript's statements that
exactly one heap is nonempty, that clearing both removes one person, and
that each iteration pops one person should all be corrected. The final
$O(n\log n)$ bound remains valid.

\subsection{The intra-group routine uses a different debtor rule}

In \texttt{\_settle}, debtors have negative balances, but
\texttt{max(debtors, key=debtors.get)} chooses the debtor closest to zero,
not the largest debtor. It also rebuilds dictionaries and scans balances
each iteration rather than using heaps. Describe the implemented rule or
change it in a separate implementation task. This discrepancy does not
invalidate feasibility or optimality within indecomposable exact blocks.

\section{Complexity and heuristic claims}

\begin{enumerate}
\item \textbf{Correct the submask identity.} The exact transition count is
\[
 \sum_{w=1}^{n}\binom nw2^{w-1}=\frac{3^n-1}{2},
\]
not $3^n/2$. The asymptotic bound is unchanged.

\item \textbf{Make space bounds consistent.} Enumerating sums by scanning
all bits takes $O(n2^n)$ time but storing integer zero-sum masks takes
$O(2^n)$ space in the usual word model. Section 4's enumeration paragraph
and its summary table currently disagree. Mention the computational model
if bit lengths of masks or monetary amounts are included.

\item \textbf{Correct intra-group cost.} The implemented scan-based
\texttt{\_settle} takes $O(k^2)$ time for a group of size $k$.
Across groups this is $O(\sum_j k_j^2)\subseteq O(n^2)$, rather than
the table's $O(nc)$ bound. A remainder settled by the same routine also
has quadratic worst-case cost. These corrections do not change the
overall $O(3^n)$ exact bound. Materializing reconstructed groups stores
$O(n)$ person references, not merely $O(c)$.

\item \textbf{Resolve the sorting contradiction.} The heuristic section
lists sorting as $O(n^2 2^n)$ but calls $O(n2^n)$ dominant. With $Z$
stored subsets, sorting by cached cardinality has the usual bound
$O(Z\log Z)\subseteq O(n2^n)$. Moreover, the implementation enumerates
by increasing cardinality already, so this sort is redundant. Storing
explicit subsets uses $O(n2^n)$ space in the worst case.

\item \textbf{Replace the impossible heuristic example.} A size-three
subset cannot be selected before any disjoint size-two subsets in this
smallest-first algorithm: the size-two subsets are processed first.
Use the existing repository fixture instead. Its balances are
\[
 (b_A,b_B,b_C,b_D,b_E,b_F,b_G,b_H,b_I)
 =(4,1,-9,-12,5,-6,2,-6,21).
\]
An optimal three-block partition is
$\{A,G,F\}$, $\{B,E,H\}$, $\{C,D,I\}$.
Each block sums to zero, giving six payments. There are no opposite
pairs, so every block has size at least three and three blocks is optimal.
The ordered heuristic first selects $\{A,E,C\}$; the six remaining
balances have no zero-sum triple and no zero-sum pair, so they cannot
split into two blocks. This yields two groups and seven payments,
as asserted by the repository test.

\item \textbf{Qualify the greedy gap.} ``Arbitrarily large'' should specify
an additive gap, with a proved family if retained. Since nonzero people
must all be incident to payments, $\mathrm{OPT}\geq\lceil n/2\rceil$,
whereas greedy uses at most $n-1$. Its ratio is therefore less than two
for nonempty instances, even though an additive gap may grow with $n$.
The informal bridge description is not a proof.

\item \textbf{Benchmark the threshold claim.} An exponential operation
count alone does not justify ``feasible in seconds.'' At $n=18$ the
transition count is about $1.94\times10^8$, with Python-level loop work.
Provide measured runtimes, hardware, Python version, and representative
input distributions before claiming an appropriate latency. The absence
of a hard heuristic size limit does not make exponential enumeration
practical for arbitrary $n$.
\end{enumerate}

\section{Input model, edge cases, and presentation}

\begin{itemize}
\item Model participants as a nonempty ordered list of distinct identifiers,
matching CSV validation. A multiset does not define ``first'' participants,
and repeated identifiers would require occurrence-indexed shares and a
different balance formula. Clarify that direct Python callers must meet
the input preconditions: \texttt{net\_balances} does not call the full
CSV validator.
\item A single participant gives no effect only when that participant is
the payer. If payer $P$ pays $A$ for sole participant $Q\ne P$, their
balances change by $+A$ and $-A$.
\item Rounding permits a share of $\lceil A/m\rceil$, so
$s_p\leq A/m$ is false. A payer-participant has a nonnegative net
contribution from that expense, which need not be positive: with one
cent split between two people and the payer first, the payer's share is
one cent and their net is zero. Include $A=0$ explicitly.
\item In the positivity invariant use $\min(-b_d,b_c)$ or
$\min(d_d,c_c)$ consistently. If $d_d$ already denotes positive debt,
negating it makes the asserted positive amount negative.
\item The dispatch display omits the explicitly supported
\texttt{heuristic} method. Explicit \texttt{exact1} bypasses the threshold;
the threshold governs only \texttt{auto}. Reflect this in the comparison
table's ``size limit'' row.
\item Qualify the abstract's global-minimum promise: it applies to exact
mode and to auto only while auto selects exact. Feasibility proofs for
approximate methods should be distinguished from optimality proofs.
\item Use \verb|\label| and \verb|\ref| throughout. For example, the
zero-sum invariant is Theorem 2.1, not simply Theorem 2 under the current
numbering. The source-reference table's ``Lines'' column actually contains
section references; rename it and use real cross-references. Remove
unused packages where convenient.
\end{itemize}

\section{Research literature and recent related work}

\subsection{Scope and search limitations}

A targeted literature search on 14 November 2026 used the terms
\emph{debts' clearing problem}, \emph{minimum number of transactions},
\emph{optimal account balancing}, and \emph{zero-sum set packing},
including searches for 2023--2026 publications and the exact phrase
\emph{Maximum Zero-Sum Partition}. The 2024 paper by Fertin and
coauthors is directly relevant, alongside earlier debt-clearing research.
Recent work with different objectives is distinguished below.
This is a selective review, not an exhaustive systematic
survey or a claim that no newer exact algorithm exists. Dates below
refer to the publications, rather than search-engine crawl dates.

\subsection{Research on the same minimum-transaction objective}

Fertin, Fontaine, Jean, and Vialette's 2024 journal
paper~\cite{fertin2024}, \emph{The Maximum Zero-Sum Partition problem},
is a central recent reference. It studies a multiset of nonzero integers
with total sum zero, maximizing the number of zero-sum blocks. Its
motivation comes from genomics, and it studies complexity,
(in)approximability, and parameterized tractability with respect to
solution size, the number of entries of either sign, and the largest
integer. A preliminary version appeared at ICS 2022. This summary is
limited to the verified publisher abstract; full-text access was blocked
during this review, so no specific theorem-level runtime or hardness
classification is attributed here.

For the present application, the equivalence is exact: map each nonzero
person balance to a distinct occurrence in the multiset, retaining person
identities even when values coincide. An MZSP optimum of $c$ blocks gives
the settlement optimum $n-c$. Conversely, connected components of an
optimal settlement yield an optimal zero-sum partition, by the graph
argument above. The manuscript should therefore explicitly identify MZSP
as its underlying combinatorial problem and cite this paper before the
more distant financial-clearing literature.

There are two useful consequences for interpreting related results.
First, the partition count $c$ is not the payment budget $K$: a decision
question about at most $K$ payments corresponds to at least $n-K$ blocks.
Parameterized results must preserve this distinction. Second, equivalent
exact objectives do not have identical approximation ratios. If a
partition method returns $\widehat c$ blocks and greedy settlement within
each block produces $T$ payments, then
\[
 0\leq T-(n-c)\leq c-\widehat c.
\]
Thus additive partition loss bounds additive payment loss. A multiplicative
guarantee for maximizing $c$ must instead be translated through $n-c$;
it should not simply be reused as the same ratio for payment minimization.
These are consequences of the equivalence, not quotations of the paper's
theorems.

P\u{a}tca\c{s}'s 2011 paper~\cite{patcas2011} studies the same objective
of clearing debts with as few transfers as possible. It relates this
objective to a maximum-cardinality zero-sum partition and develops
subset-based methods for changing debt graphs. This provides prior
research context for the manuscript's graph model, $n-c$ characterization,
and exponential DP. Its dynamic setting also motivates maintaining net
balances when expenses change. ``Dynamic'' here concerns updates to the
input graph, and should not be confused with the economic scheduling
models discussed below.

Yao's 2017 senior thesis~\cite{yao2017} investigates the connection to
zero-sum set packing, complexity, approximation methods, and restrictions
on permissible payment edges. It is a useful source for the distinction
between unrestricted settlement and graph-constrained variants. For this
project, its main methodological relevance is that the zero-sum packing
formulation has existing literature; presenting it with attribution is
preferable to suggesting that it is new. The source is a university thesis,
not a journal article.

P\u{a}tca\c{s} and Bartha's 2019 journal paper~\cite{patcas2019}
proposes an evolutionary approach for the minimum-transaction problem,
which it identifies as strongly NP-hard. It is directly relevant when
comparing methods intended for larger groups. Such a method is a candidate
for experimental comparison with the repository's subset-packing and heap
heuristics, not a replacement for the exact optimality guarantee. Any
empirical comparison should report repeated runs and computation budgets
for stochastic methods, alongside payment counts and verified feasibility.

These references support adding a complexity discussion: the decision
version asks whether a settlement with at most $K$ transfers exists;
the optimization problem asks for the smallest count. Avoid calling an
optimization problem ``NP-complete'' without specifying its decision
version. NP-hardness is compatible with the faster exponential DP already
described in this review and does not establish that $O(3^n)$ is necessary.

\subsection{Recent research with related, but different, objectives}

Chen, Baron, and Chen~\cite{chen2025} study a dynamic interbank system
where obligations arrive over time. Their objective combines clearing
setup costs, costs depending on net debt, and liquidity costs from leaving
debts uncleared. They use state reduction and dynamic programming to
analyze clearing policies. The paper appeared online in June 2024 and in
the April 2025 journal issue. It is relevant if the application later
chooses \emph{when} to settle, but its policy problem does not supply an
exact minimum-transfer algorithm for the current static balance vector.

Buchman and coauthors' \emph{Cycles Protocol}~\cite{buchman2025} is a
2025 arXiv preprint describing multilateral clearing using obligation
graphs, multiple liquidity sources, and privacy-preserving settlement.
Its emphasis is discharging debt with less liquidity rather than minimizing
the number of positive transfers after unrestricted netting. It is useful
context for a future settlement platform, but does not establish improved
complexity or optimality guarantees for the manuscript's objective. The
arXiv submission is dated July 2025; the manuscript itself displays a
November 2024 version date. The citation below identifies the arXiv
version, without asserting peer-reviewed publication.

Both recent works illustrate why the review should state its objective
precisely. Minimizing transaction count, minimizing liquidity requirements,
and minimizing the cost of delaying settlement are distinct problems.
Results should only be transferred between them after establishing the
necessary reduction and assumptions.

\subsection{Practical exact-algorithm reference and implications}

The 2025 official AtCoder editorial~\cite{atcoder2025} explicitly gives
an $O(n2^n)$-time, $O(2^n)$-space method for the equivalent zero-balance
partitioning problem. It is an algorithmic exposition, not research
literature, and should be labeled accordingly. Its inclusion verifies
that the faster exact method has a public prior description; it is not
a claim about who first discovered that method or about the best known
exact complexity in all computational models.

The methodology should therefore cite the 2024 MZSP paper and prior
debt-clearing research,
include the faster exact DP as a concrete alternative to its implementation,
and retain heuristic methods as explicitly approximate choices. Recent
liquidity and scheduling research belongs in related work rather than
in the table comparing algorithms for the identical static objective.

\section{Recommended revision and validation}

First discuss the faster $O(n2^n)$ exact DP and distinguish it from the
current implementation. Repair the graph proof and theorem quantifiers,
then establish the
DP recurrence and the zero-sum reconstruction invariant. Next align both
greedy descriptions with their implementations and reconcile the complexity
tables. Finally replace unsupported examples and timing claims with
checkable evidence.

For a future implementation validation pass, useful cases include the empty
instance, a sole participant different from the payer, one-cent rounding,
the two-person heap-sign example, the nine-person heuristic fixture, and
generic reconstruction from a mask that requires skipping its lowest bit.
For small zero-sum instances, compare exact output counts against an
independent exhaustive partition oracle, as well as checking that applying
every payment returns all balances to zero. These checks would support
the repaired claims; this review does not claim to have run them.

\begin{thebibliography}{9}

\bibitem{fertin2024}
Guillaume Fertin, Oscar Fontaine, G\'{e}raldine Jean, and St\'{e}phane Vialette.
\emph{The Maximum Zero-Sum Partition problem}.
Theoretical Computer Science, 1019:114811, 1 December 2024.
\url{https://doi.org/10.1016/j.tcs.2024.114811}.
Author-deposited record: \url{https://hal.science/hal-04901397}.

\bibitem{patcas2011}
Csaba P\u{a}tca\c{s}.
\emph{The debts' clearing problem: a new approach}.
Acta Universitatis Sapientiae, Informatica, 3(2):192--204, 2011.
\url{https://arxiv.org/abs/1111.3663}.

\bibitem{yao2017}
Y.~Yao.
\emph{Settling Debts Efficiently: Zero-Sum Set Packing}.
Senior thesis, Harvard University, 2017.
\url{https://dash.harvard.edu/bitstream/handle/1/38811480/YAO-SENIORTHESIS-2017.pdf}.

\bibitem{patcas2019}
Csaba P\u{a}tca\c{s} and Attila Bartha.
\emph{Evolutionary solving of the debts' clearing problem}.
Acta Universitatis Sapientiae, Informatica, 11(2):142--158, 2019.
\url{https://doi.org/10.2478/ausi-2019-0010}.
An earlier preprint is available as arXiv:1402.6556 (2014).

\bibitem{chen2025}
Shuzhen Chen, Opher Baron, and Ningyuan Chen.
\emph{Optimal Dynamic Clearing for Interbank Payments}.
Management Science, 71(4):2953--2974, 2025.
First published online 27 June 2024.
\url{https://doi.org/10.1287/mnsc.2023.00380}.

\bibitem{buchman2025}
Ethan Buchman, Paolo Dini, Shoaib Ahmed, Andrew Miller, and
Toma\v{z} Fleischman.
\emph{Cycles Protocol: A Peer-to-Peer Electronic Clearing System}.
arXiv preprint arXiv:2507.22309v1, 30 July 2025.
\url{https://arxiv.org/abs/2507.22309v1}.

\bibitem{atcoder2025}
spheniscine.
\emph{F---Candy Redistribution: Solution in O(N * 2\^{}N)}.
Official editorial, AtCoder Beginner Contest 432, 2025.
Programming-contest exposition.
\url{https://atcoder.jp/contests/abc432/editorial/14584?lang=en}.

\end{thebibliography}
\end{document}
